\documentclass[11pt]{article}

\usepackage[margin=1in]{geometry}
\usepackage{amsmath,amssymb,amsthm,mathtools,microtype}
\usepackage[T1]{fontenc}
\usepackage[hidelinks]{hyperref}

\newtheorem{theorem}{Theorem}[section]
\newtheorem{lemma}[theorem]{Lemma}
\newtheorem{corollary}[theorem]{Corollary}
\newtheorem{remark}[theorem]{Remark}

\newcommand{\C}{\mathbb C}
\newcommand{\Z}{\mathbb Z}
\newcommand{\N}{\mathbb N}
\newcommand{\E}{\mathbb E}
\newcommand{\CT}{\operatorname{CT}}
\newcommand{\ord}{\operatorname{ord}}
\newcommand{\conv}{\operatorname{conv}}
\newcommand{\Supp}{\operatorname{Supp}}

\title{The Gaussian Moments Conjecture:\\
Dimension Two and Sparse Minimality in Dimension Three}
\author{Roy van Rijn\\{\small Independent researcher}}
\date{24 July 2026}
\hypersetup{
  pdftitle={The Gaussian Moments Conjecture: Dimension Two and Sparse Minimality in Dimension Three},
  pdfauthor={Roy van Rijn},
  pdfsubject={The dimension threshold and sparse minimal counterexamples for the Gaussian Moments Conjecture},
  pdfkeywords={Gaussian Moments Conjecture, Gaussian moments, constant terms,
    Laurent polynomials, sparse polynomials, Groebner bases}
}

\begin{document}
\maketitle

\begin{abstract}
The Gaussian Moments Conjecture asks whether a polynomial whose positive
powers all have Gaussian integral zero must eventually remain zero after
multiplication by any fixed polynomial.  We prove the conjecture for two
real Gaussian variables.  The proof uses circular coordinates to separate
rotation from radius.  The angular part becomes a constant-term problem,
while the radial part becomes the simple functional \(U^j\mapsto j!\).
An exposed lower face of the resulting Newton diagram, the constant-term
theorem of Duistermaat and van der Kallen, and reduction modulo a prime then
isolate a nonzero lowest term.  This forces every polynomial with vanishing
pure moments to have rotational weights of one strict sign, from which the
mixed-moment conclusion is immediate.  Together with the recent
counterexamples of Long in every dimension at least three, this determines
the exact dimensional range of the conjecture.  We then study minimal
failure in the first failing dimension.  An exhaustive exact search over
all \(59{,}535\) supports having degree at most four and at most four
monomials shows that every pure-moment-null polynomial in this range still
satisfies the conjectural mixed-moment conclusion.  Thus Long's five-term
quartic is Pareto-minimal with respect to degree and circular-coordinate
support size.  The sparse theorem is computer-assisted: saturated moment
ideals through order twenty-four are computed over \(\mathbb Q\), and the
three genuine surviving families are proved harmless at all orders.
\end{abstract}

\section{Introduction}

Let \(X_1,\ldots,X_d\) be independent standard real Gaussian random
variables.  The Gaussian Moments Conjecture, introduced by Derksen, van den
Essen, and Zhao~\cite{DVZ}, is the following statement.

\medskip
\noindent
\textbf{Gaussian Moments Conjecture \(\operatorname{GMC}(d)\).}
If \(P\in\C[X_1,\ldots,X_d]\) satisfies
\[
  \E(P^m)=0\qquad\text{for every }m\geq 1,
\]
then, for every \(Q\in\C[X_1,\ldots,X_d]\),
\[
  \E(QP^m)=0\qquad\text{for all sufficiently large }m.
\]

\medskip
The conjecture was motivated in part by its connection with the Jacobian
Conjecture~\cite{DVZ}.  Long recently found explicit counterexamples to
\(\operatorname{GMC}(d)\) for every \(d\geq3\)~\cite{Long}.  The
two-variable case was therefore the last undecided dimension.

\begin{theorem}\label{thm:main}
The Gaussian Moments Conjecture holds in two real variables.
\end{theorem}

Long's three-variable witness, in circular coordinates \(W,Z\) and an
independent real Gaussian \(T\), is
\begin{equation}\label{eq:long}
 P_{\mathrm L}
 =W+WZ-T^2-\frac32ZT^2-\frac12Z^2T^2.
\end{equation}
It has degree four and five monomials, all its pure moments vanish, and
\(\E(ZP_{\mathrm L}^m)=m!\) for every \(m\geq1\)~\cite{Long}.  Our second
main result places this witness at the first corner of the sparse
degree--support plane.

\begin{theorem}[Sparse Pareto minimality]\label{thm:sparse}
Let \(W,Z,T\) be Gaussian Witt coordinates, so that
\[
 \E(W^iZ^jT^k)=
 \begin{cases}
   i!(k-1)!!,&i=j\text{ and }k\text{ is even},\\
   0,&\text{otherwise}.
 \end{cases}
\]
Suppose that \(P\in\C[W,Z,T]\) has total degree at most four and at most four
nonzero monomials.  If \(\E(P^m)=0\) for every \(m\geq1\), then, for every
\(Q\in\C[W,Z,T]\),
\[
 \E(QP^m)=0
 \qquad\text{for all sufficiently large }m.
\]
Consequently no GMC counterexample has both degree at most four and support
size at most four in a Gaussian Witt frame.
\end{theorem}

In particular, \eqref{eq:long} is minimal for the product partial order on
\((\deg P,\#\Supp(P))\).  The theorem does not separately prove that degree
four is globally minimal: a denser cubic is not excluded.  Nor does it
exclude a four-term counterexample of degree greater than four.

The proof is short once the right coordinates are chosen.  In two
dimensions, rotation has a single integer weight and there is a single
radial variable.  Thus every polynomial can be displayed schematically as
\[
  P=\sum_k \underbrace{T^k}_{\text{rotation}}\,
                  \underbrace{B_k(U)}_{\text{radius}}.
\]
Gaussian integration first takes the constant term in \(T\), then replaces
\(U^j\) by \(j!\).  If both positive and negative rotational weights occur,
the lowest radial orders have a supporting line above weight zero.  The
polynomial carried by that lower face has a nonzero constant term in some
power, by a theorem of Duistermaat and van der Kallen~\cite{DK}.  Raising
that power to a suitable prime and reducing modulo the prime makes the
lowest term impossible to cancel after the factorial substitution.  This
contradicts the assumed moment vanishing.

The argument is conceptual rather than computational.  In particular, it
contains the previously studied three-level and star-shaped supports
without needing separate case analyses.

\section{Circular coordinates}

Let \(X,Y\) be independent standard real Gaussian variables.  Set
\[
  Z=\frac{X+iY}{\sqrt2},\qquad
  W=\frac{X-iY}{\sqrt2}.
\]
Wick's formula gives
\begin{equation}\label{eq:wick}
  \E(Z^aW^b)=
  \begin{cases}
    a!,&a=b,\\
    0,&a\neq b.
  \end{cases}
\end{equation}
Introduce a radial variable \(U\) and a Laurent variable \(T\) through the
injective substitution
\[
  Z\longmapsto T,\qquad W\longmapsto UT^{-1}.
\]
Thus \(Z^aW^b\) becomes \(T^{a-b}U^b\).  Every polynomial \(P(Z,W)\)
therefore has a unique expression
\begin{equation}\label{eq:weight-decomposition}
  P=\sum_{k\in S}T^kB_k(U),
\end{equation}
where \(S\subset\Z\) is finite and every \(B_k\in\C[U]\) is nonzero.  The
integer \(k\) is the rotational weight.

Define the factorial functional
\[
  \mathcal L:\C[U]\longrightarrow\C,\qquad
  \mathcal L(U^j)=j!.
\]
Equation~\eqref{eq:wick} is exactly the identity
\begin{equation}\label{eq:expectation-factorization}
  \E(F)=\mathcal L\bigl(\CT_T F\bigr)
  \qquad(F\in\C[Z,W]).
\end{equation}
This is the feature special to two variables: there is one angular
constant-term operation and one radial factorial.

\section{The exposed lower face}

For each \(k\in S\), write
\[
  \nu_k=\ord_U B_k,\qquad
  b_k=[U^{\nu_k}]B_k\neq0.
\]
The points \((k,\nu_k)\) record the smallest radial degree at each
rotational weight.  We now extract the lower face visible from weight zero.

\begin{lemma}[Lower face]\label{lem:face}
Suppose \(0\in\conv(S)\).  There are rational numbers \(\rho,\theta\), a
subset
\[
  S_0=\{k\in S:\nu_k=\rho+\theta k\},
\]
and an integer \(r\geq1\) with the following properties:
\[
  \nu_k\geq\rho+\theta k\quad(k\in S),\qquad
  0\in\conv(S_0),
\]
and, for
\[
  P_0(T)=\sum_{k\in S_0}b_kT^k,\qquad
  F_m(U)=\CT_T(P^m),
\]
there are \(n\in\Z_{\geq0}\) and \(c\in\C^\times\) such that
\begin{align}
  &\text{every term of }F_m\text{ has degree at least }\rho m,
       \label{eq:global-bound}\\
  &n=\rho r,\qquad F_r(U)=cU^n+\text{terms of degree greater than }n.
       \label{eq:leading}
\end{align}
\end{lemma}

\begin{proof}
Minimize
\[
  \sum_{k\in S}\lambda_k\nu_k
\]
subject to
\[
  \lambda_k\geq0,\qquad
  \sum_k\lambda_k=1,\qquad
  \sum_k k\lambda_k=0.
\]
The feasible set is nonempty because \(0\in\conv(S)\).  Let the minimum be
\(\rho\).  Elementary linear-programming duality supplies
\(\theta\in\mathbb Q\) such that
\[
  \nu_k\geq\rho+\theta k.
\]
An optimal solution is supported on the equality set \(S_0\), so
\(0\in\conv(S_0)\).

We use the constant-term theorem of Duistermaat and van der
Kallen~\cite{DK}: if every positive power of a complex Laurent polynomial
has zero constant term, then zero does not lie in the convex hull of its
support.  Applied contrapositively to \(P_0\), it gives an \(r\geq1\) for
which
\[
  c:=\CT_T(P_0^r)\neq0.
\]

Consider a weight-zero product of \(m\) terms selected from
\eqref{eq:weight-decomposition}.  If the selected weights are
\(k_1,\ldots,k_m\), then its radial degree is at least
\[
  \sum_{a=1}^m\nu_{k_a}
  \geq \rho m+\theta\sum_{a=1}^m k_a
  =\rho m.
\]
This proves~\eqref{eq:global-bound}.  When \(m=r\), equality can occur only
by taking the lowest coefficient from weights on \(S_0\), and the sum of
all equality terms is \(c\).  Hence the lowest term of \(F_r\) is
\(cU^{\rho r}\).  Its exponent is necessarily a nonnegative integer; call
it \(n\).
\end{proof}

\section{Prime isolation}

The next lemma is the arithmetic step.  Its role is simple: after a prime
dilation, Frobenius preserves only \(p\)-th powers, while the factorial
functional kills every radial level above the lowest one.

\begin{lemma}[Prime--factorial isolation]\label{lem:prime}
Let \(P\) and \(F_m\) be as in Lemma~\ref{lem:face}.  If
\[
  \mathcal L(F_m)=0\qquad(m\geq1),
\]
then the conclusion~\eqref{eq:leading} is impossible.
\end{lemma}

\begin{proof}
Let \(R\subset\C\) be the \(\Z\)-algebra generated by the coefficients of
\(P\) and by \(c^{-1}\).  It is a finitely generated domain of
characteristic zero.  For all but finitely many rational primes \(p\),
there is a homomorphism from \(R\) to a field of characteristic \(p\).
Indeed, the finite-type morphism
\(\operatorname{Spec}R\to\operatorname{Spec}\Z\) is dominant, and its image
is constructible by Chevalley's theorem~\cite{StacksChevalley}.  Since it
contains the generic point, it contains a nonempty open
subset~\cite{StacksGeneric}.  Because \(c\) is invertible in \(R\), its
image remains nonzero under every such homomorphism.

Fix one of these primes.  Before reducing modulo \(p\), work in the
characteristic-zero domain \(R\).  Lemma~\ref{lem:face} says that every term
of \(F_{rp}\) has degree at least \(\rho rp=np\).  Hence
\(\mathcal L(F_{rp})=0\) can be written in \(R\) as
\[
  0=(np)!\sum_{d\geq np}[U^d]F_{rp}\,\frac{d!}{(np)!}.
\]
Since \(R\) is a domain of characteristic zero, \((np)!\neq0\), so it may
be cancelled to give
\begin{equation}\label{eq:normalized}
  \sum_{d\geq np}[U^d]F_{rp}\,\frac{d!}{(np)!}=0.
\end{equation}
We now apply a homomorphism \(R\to\kappa\), where \(\kappa\) is a field of
characteristic \(p\), and continue to use the same notation for the reduced
polynomials.

In characteristic \(p\), Frobenius and constant-term extraction give
\begin{equation}\label{eq:frobenius}
  F_{rp}
  =\CT_T\bigl((P^r)^p\bigr)
  =\bigl(\CT_T(P^r)\bigr)^p
  =F_r^p.
\end{equation}
Indeed, if \(H=\sum_i h_i(U)T^i\), then
\(H^p=\sum_i h_i(U)^pT^{pi}\) in characteristic \(p\); because \(p>0\),
the equality \(pi=0\) holds exactly when \(i=0\).  Thus
\(\CT_T(H^p)=\CT_T(H)^p\).
By~\eqref{eq:leading}, the coefficient of \(U^{np}\) in
\eqref{eq:frobenius} is \(c^p\).  A term whose degree is not divisible by
\(p\) has zero coefficient modulo \(p\).  A remaining term has degree
\(jp\) for some \(j>n\), and its multiplier in~\eqref{eq:normalized} is
\[
  \frac{(jp)!}{(np)!}.
\]
This integer is divisible by \(p\), since its product contains
\((n+1)p\).  Consequently, reducing~\eqref{eq:normalized} modulo \(p\)
leaves only
\[
  c^p=0,
\]
contrary to the choice of the reduction.
\end{proof}

\section{Proof of the theorem}

\begin{proof}[Proof of Theorem~\ref{thm:main}]
Write \(P\) as in~\eqref{eq:weight-decomposition} and suppose that
\(\E(P^m)=0\) for every \(m\geq1\).  By
\eqref{eq:expectation-factorization}, this says
\[
  \mathcal L(F_m)=0\qquad(m\geq1).
\]
If \(0\in\conv(S)\), Lemma~\ref{lem:face} produces a nonzero lowest
coefficient and Lemma~\ref{lem:prime} gives a contradiction.  Therefore
\[
  0\notin\conv(S).
\]
Since \(S\subset\Z\), either every weight of \(P\) is strictly positive or
every weight is strictly negative.

Now fix a polynomial \(Q\).  Its rotational weights form a bounded finite
set.  Every weight of \(P^m\), on the other hand, tends uniformly toward
the same strict side as \(m\to\infty\).  Hence \(QP^m\) has no
weight-zero term once \(m\) is sufficiently large:
\[
  \CT_T(QP^m)=0.
\]
Equation~\eqref{eq:expectation-factorization} then gives
\(\E(QP^m)=0\), as required.
\end{proof}

\begin{corollary}\label{cor:classification}
The Gaussian Moments Conjecture holds in dimension \(d\) exactly when
\(d\leq2\).
\end{corollary}

\begin{proof}
The case \(d=1\) is known~\cite[Proposition~4.2]{DVZ}.
Theorem~\ref{thm:main} proves the case \(d=2\).
Long's examples disprove the conjecture for every \(d\geq3\)~\cite{Long}.
\end{proof}

\section{Sparse minimality in the first failing dimension}
\label{sec:sparse}

We prove Theorem~\ref{thm:sparse}.  The finite part is an exhaustive exact
coefficient-scheme calculation; the all-order interpretation of its
survivors is conceptual.

\subsection{Moment ideals}

Let
\[
 \mathcal M_4
 =\{W^iZ^jT^k:i,j,k\geq0,\ i+j+k\leq4\}.
\]
This set has \(35\) elements.  Fix a support \(S\subset\mathcal M_4\), write
\[
 P_S=\sum_{s\in S}c_ss,
\]
and work on the coefficient torus \(\prod_{s\in S}c_s\neq0\).  For
\(s=(i_s,j_s,k_s)\) and a multiplicity vector
\(\nu=(\nu_s)_{s\in S}\in\mathbb N^S\), put
\[
 |\nu|=\sum_s\nu_s,\quad
 I(\nu)=\sum_s\nu_si_s,\quad
 J(\nu)=\sum_s\nu_sj_s,\quad
 K(\nu)=\sum_s\nu_sk_s.
\]
Wick contraction gives the exact moment polynomial
\begin{equation}\label{eq:sparse-moment}
 \E(P_S^m)
 =
 \sum_{\substack{|\nu|=m,\ I(\nu)=J(\nu)\\K(\nu)\ {\rm even}}}
 \frac{m!}{\prod_s\nu_s!}\,
 I(\nu)!\,(K(\nu)-1)!!\prod_sc_s^{\nu_s}.
\end{equation}
Thus the truncated full-support coefficient scheme is defined by
\begin{equation}\label{eq:saturated-ideal}
 I_{S,M}
 =
 \left\langle\E(P_S),\ldots,\E(P_S^M)\right\rangle:
 \left(\prod_{s\in S}c_s\right)^\infty.
\end{equation}
The equations are homogeneous, so one coefficient may be normalized to
one before saturation.

There are
\[
 \sum_{r=1}^4\binom{35}{r}=59{,}535
\]
raw supports.  We identify a support with its reflection under
\(W\leftrightarrow Z\).  Two elementary filters precede the Groebner
calculation.  If some moment in \eqref{eq:sparse-moment} has exactly one
coefficient monomial, the coefficient torus contains no solution.  If no
contraction occurs, all charges \(i_s-j_s\) have one strict sign and the
mixed-moment conclusion follows by bounded weight.

For completeness, the latter test is finite in this range.  A zero charge
contracts in one or two copies, according to its \(T\)-parity.  Opposite
nonzero charges \(q>0>r\), with \(|q|,|r|\leq4\), have a primitive
zero-sum relation of length
\[
 \frac{q+|r|}{\gcd(q,|r|)}\leq8;
\]
doubling corrects odd \(T\)-parity.  Hence every support which is not
strictly one-sided has a contraction by order sixteen.

\subsection{Exact enumeration}

After the preceding filters, the four-term stratum contains \(5{,}718\)
coefficient-torus ideals requiring elimination.  Exact Groebner bases over
\(\mathbb Q\) for \(I_{S,20}\) leave precisely the supports
\[
\begin{split}
 &\{Z,ZT^2,WZ^2,W^3T\},\\
 &\{ZT,Z^2,WT,W^2\},\\
 &\{ZT,ZT^3,WZ^2T,W^3\}.
\end{split}
\]
The first and third acquire the unit ideal after adjoining moments through
order twenty-four.  The same exact search in the one-, two-, and three-term
strata leaves no active one- or two-term family and two three-term
families.  Including the four-term survivor, the genuine active null
families are
\begin{align}
 P_{\rm lin}&=T+aZ+bW,&1+2ab&=0,                 \label{eq:null-linear}\\
 P_{\rm sh}&=Z+aWT-\frac12a^2W^3,                \label{eq:null-shear}\\
 P_{\rm pr}&=ZT+aZ^2+\frac1{2a^2}WT
                    -\frac1{4a^3}W^2,\quad a\neq0. \label{eq:null-product}
\end{align}
The saturated lexicographic bases are triangular.  With the first
coefficient normalized to one, their nontrivial equations are respectively
\[
 1+2c_1c_2,\qquad
 2c_2+c_1^2,\qquad
 2c_1^2c_2-1,\quad c_3+c_1c_2^2.
\]
These formulas both classify the coefficient schemes and make the
survivors easy to interpret.

\subsection{All-order interpretation of the survivors}

The family \eqref{eq:null-linear} is an isotropic linear form.  A complex
orthogonal change takes it to a scalar multiple of \(Z\).  A fixed
polynomial \(Q\) has bounded \(W\)-degree, so
\(\E(QZ^m)=0\) for all sufficiently large \(m\).

For \eqref{eq:null-shear}, define
\[
 D=W\partial_T-T\partial_Z.
\]
Then \(D(2WZ+T^2)=0\), \(D(W)=0\), and
\[
 P_{\rm sh}=\exp(-aWD)(Z).
\]
Gaussian integration by parts gives \(\E(DF)=0\) for every polynomial
\(F\), and therefore also \(\E(WDF)=0\).  Since \(WD\) is locally nilpotent,
\(\phi=\exp(-aWD)\) is a polynomial automorphism preserving the Gaussian
functional.  For fixed \(Q\),
\[
 \E(QP_{\rm sh}^m)
 =\E\bigl(\phi^{-1}(Q)Z^m\bigr),
\]
which vanishes eventually by the same bounded-weight argument.

Finally,
\begin{equation}\label{eq:product-factorization}
 P_{\rm pr}
 =
 \left(Z+\frac1{2a^2}W\right)
 \left(T+aZ-\frac1{2a}W\right).
\end{equation}
The second factor is isotropic and orthogonal to the first.  After an
orthogonal Witt change, \eqref{eq:product-factorization} is a scalar
multiple of \(TZ\).  Every \(P_{\rm pr}^m\) then has isotropic weight \(m\),
which a fixed \(Q\) can balance only finitely often.  All three families
satisfy GMC.  Together with the strictly one-sided supports, this proves
Theorem~\ref{thm:sparse}.

\subsection{Reproducibility and scope}

The enumeration is implemented with exact integer Wick coefficients and
Singular's rational Groebner bases.  The certificate script enumerates
supports rather than reading a stored survivor list, repeats every
saturation over \(\mathbb Q\), verifies the three displayed triangular
bases, and checks \eqref{eq:product-factorization}.  It is run from the
repository root by
\begin{verbatim}
.venv/bin/python scripts/verify_three_real_gmc_sparse_pareto.py
\end{verbatim}
The complete replay takes several minutes on a laptop.  Modular runs are
used only as optional discovery filters and are not part of the exact
certificate.

\section{What the argument says}

The proof gives a stronger structural statement than the mixed-moment
conclusion alone:
\[
  \E(P^m)=0\ \text{for all }m\geq1
  \quad\Longrightarrow\quad
  \text{all rotational weights of \(P\) have one strict sign.}
\]
Thus vanishing in two variables is not caused by a complicated balance
among positive, zero, and negative angular modes.  Such a balance always
has a lowest radial face, and a prime moment detects it.

This also explains why support-by-support calculations quickly appeared to
show rigidity.  A three-level support, a star support, or a tree of
primitive weight relations is only a finite presentation of the same lower
face.  Leaf removal is unnecessary: the supporting line exposes all
relevant leaves at once.

The proof is genuinely two-dimensional.  With one circular pair there is
one weight coordinate and one radial order, and the latter is measured by
the single factorial \(j!\).  In higher dimension there are several radial
directions, so the one-dimensional isolation step no longer applies.
Corollary~\ref{cor:classification} shows that this is not merely a
technical gap: counterexamples already occur in dimension three.

\section{Formal verification}

A companion Lean~4 development has been prepared as supplementary material
for the two-variable theorem.  It mechanically checks the
factorial-quotient divisibility, prime-dilated factorial isolation,
commutation of Frobenius with angular constant-term extraction, the final
prime contradiction, and the eventual vanishing argument for both signs of
the rotational support.  The development contains no \texttt{sorry}
declarations.  The computer-assisted sparse theorem is certified separately
by the exact rational script described in Section~\ref{sec:sparse}.

The formalization exposes three external inputs as named axioms: the
Duistermaat--van der Kallen constant-term theorem, rational supporting-face
extraction, and the existence of good finite-characteristic specializations
of finite-type \(\Z\)-domains.  Thus the artifact gives a machine-checked
verification of the arithmetic and constant-term core of the proof, with
its remaining trust boundary stated explicitly.

\section*{Acknowledgments and declaration of generative AI use}

The author used OpenAI's ChatGPT and Codex extensively during the research
process, including for conjecture formation, symbolic experimentation,
proof exploration, consistency checks, and preparation of this manuscript.
The lower-face and prime-isolation argument emerged from that collaborative
process.  The author reviewed the mathematical claims and assumes full
responsibility for the contents.  The proof of Theorem~\ref{thm:main} is
noncomputational.  Theorem~\ref{thm:sparse} uses the exhaustive exact
coefficient-scheme calculation stated in the text and supplied with the
paper.

\bibliographystyle{plain}
\bibliography{references}

\end{document}
